\par\Needspace{7\baselineskip}\SourceLeaf{405}
\begin{center}\textsc{Gauss} to \textsc{Olbers}. Göttingen, 30 October 1825.\end{center}
\indent\mbox{. . . . . . . . . .} I must have expressed myself unclearly in my letter about my reservations concerning the arrangement of my projected work on Higher Geodesy. The question troubling me is not whether to include my view of the radii of curvature. It is whether to include the very general investigations that keep growing in my hands, into curved surfaces, the figures formed on them, the nature of shortest and non-shortest lines, and a host of other subjects I cannot list here because their concepts are not yet current and do not even have names --- whether to include all of this, especially since every day I find more subjects to which I must go back. To my knowledge they have not previously been considered from the standpoint needed to connect the whole. One of them (a minor example) was even the theory of radii of curvature in Plano.

I might also have mentioned the theory of the area of figures in general, which I likewise have considered for 30 years and more from a standpoint I had until now thought new. I was partly mistaken about that. Only recently did I become acquainted with a paper by \textsc{Meister} (a very ingenious mind, in my opinion) in the first volume of Novi Commentarii Gotting., where the subject is considered in almost the same way and is developed very beautifully.

Yet this excellent paper is almost entirely unknown to mathematicians. It would not be enough for me to refer to it, since it gives only the first basic concepts. In the analytical treatment of geometrical relations, one has long understood lines of positive and negative value quite well and recognized that a certain sense (sens, direction) is always taken as the basis, explicitly or implicitly. The lines are regarded as increasing in that direction, etc. Yet when surface areas (areas) are expressed by formulas, a %
